\section{维数正规化下含“重夸克”的有效 QCD 的重正化}

我们来考虑上述含“重夸克”有效理论的重正化。
我们将证明这样的理论可以在微扰论所有阶进行微扰重正化，
首先在不出现幂次发散的维数正规化中进行。基本论证与 HQET 可重正化的全阶证明相同，
后者由文献~\cite{Bagan:1993zv} 首先给出。

QCD 可重正化的标准证明选取协变规范 $\partial_\mu A^\mu = 0$，并在泛函积分中以规范参数 $\xi$ 实现。为了在规范固定之后恢复定域理论，必须引入色八重态鬼场 $c^a$ 和 $\bar c^a$。所得理论具有一种称为 BRST 对称性的剩余整体对称性。来自这一对称性的 Ward-Takahashi 恒等式，是证明理论在微扰论所有阶可重正化的关键要素，其含义如下：把所有无穷大吸收进夸克（$Z_1$）、胶子（$Z_3$）和鬼场（$Z_c$）的波函数重正化因子，以及夸克质量（$Z_m$）、规范耦合（$Z_g$）和规范固定参数（$Z_\xi$）的乘法重正化之后，重正化基本场的格林函数是 UV 有限的。

在 HQET 中，拉氏量的形式与式~(\ref{Eq:lag}) 类似，只是矢量 $n^\mu$ 是类时的而非类空的；已经证明，在重夸克质量的领头阶，格林函数可以在微扰论所有阶重正化~\cite{Bagan:1993zv}，其中用到 HQET 拉氏量的 BRST 不变性。在 HQET 中，重夸克不与重反夸克耦合，因此不会出现重夸克圈图。这意味着不含外部重夸克腿的格林函数可以像在 QCD 中那样被重正化。只有那些涉及外部重夸克腿的格林函数才需要单独考虑，它们给出一个新的波函数重正化常数（$Z_Q$）。这些格林函数包括重夸克两点函数、重夸克-胶子顶点，以及带有一个来自重夸克场 BRST 变换的复合算符插入的重夸克-鬼场顶点。借助 Ward-Takahashi 恒等式，在 HQET 拉氏量给定的有限个重正化因子下，它们全部可以证明是 UV 有限的~\cite{Bagan:1993zv}。

注意，上述论证的有效性与 HQET 拉氏量中的重夸克场是由类时还是类空矢量来定义无关。
因此，文献~\cite{Bagan:1993zv} 关于 HQET 可重正化的论断，对我们的辅助重夸克拉氏量（式~(\ref{Eq:lag}) 中 $n^\mu=(0,0,0,1)$）同样成立。然而，在光锥矢量的情形下，由于光锥奇异性会出现新的发散。

该有效理论中的定域复合算符可以用标准方法重正化。特别地，轻重复合算符，例如
\begin{equation}
             \overline{j}(x) = \overline{\psi}(x)\Gamma Q(x),
\end{equation}
会获得一个发散因子 $Z_j$（注意此时 $\overline{j}(x)$ 带有一个旋量指标）。类比于格林函数的情形，可以证明：把 $\bar{j}(x)$ 插入格林函数所产生的整体 UV 发散，在微扰论所有阶都被 $Z_j$ 减除。因此，我们可以写 $\overline{j}= Z_j \overline{j}_R$，其中 $\overline{j}_R$
是重正化的算符。
 其反常量纲的计算
已在 HQET 中完成~\cite{Politzer:1988wp,Ji:1991pr}。
对于辅助重夸克场，实际上可以显式验证：在维数正规化（取 $D=4-2\epsilon$）下的单圈水平上，
\begin{equation}\label{eq:hl}
   Z_j = 1+ \frac{\alpha_s}{2\pi\epsilon} \ ,
\end{equation}
这与 HQET 中的结果相同。
